%!TEX program = xelatex
\documentclass[dvipsnames, svgnames,a4paper,11pt]{article}
% ----------------------------------------------------
%   中山大学物理与天文学院本科实验报告模板
%   作者：Huanyu Shi，2019级
%   知乎：https://www.zhihu.com/people/za-ran-zhu-fu-liu-xing
%   Github：https://github.com/huanyushi/SYSU-SPA-Labreport-Template
%   Last update : 2023.4.10
% ----------------------------------------------------

\input{Settings} % 导入模板的相关设置
\usepackage{lipsum}
\usepackage{enumitem}
\usepackage{tabularray}  % 绘制表格时可以更加方便添加框线
\usepackage{diagbox}	% 方便给表格添加对角线
% \usepackage{subfig}
% \usepackage{caption}
\setlist[enumerate]{label=\textup{(\arabic*)}}



%---------------------------------------------------------------------
%	正文
%---------------------------------------------------------------------

\begin{document}


\begin{table}
	\renewcommand\arraystretch{1.7}
	\begin{tabularx}{\textwidth}{
		|X|X|X|X
		|X|X|X|X|}
	\hline
	\multicolumn{2}{|c|}{预习报告}&\multicolumn{2}{|c|}{实验记录}&\multicolumn{2}{|c|}{分析讨论}&\multicolumn{2}{|c|}{总成绩}\\
	\hline
	\LARGE25 & & \LARGE30 & & \LARGE25 & & \LARGE80 & \\
	\hline
	\end{tabularx}
\end{table}


\begin{table}
	\renewcommand\arraystretch{1.7}
	\begin{tabularx}{\textwidth}{|X|X|X|X|}
	\hline
	专业：& 物理学 &年级：& 2022级\\
	\hline
	姓名：& 戴鹏辉  & 学号： & 22344016 \\
	\hline
	日期：& 2024/10/22 & 教师签名：& \\
	\hline
	\end{tabularx}
\end{table}

\begin{center}
	\LARGE D1-A \quad 电阻热噪声及玻尔兹曼常量测量
\end{center}

\textbf{【实验报告注意事项】}
\begin{enumerate}
	\item 实验报告由三部分组成：
	\begin{enumerate}
		\item 预习报告：（提前一周）认真研读\underline{\textbf{实验讲义}}，弄清实验原理；实验所需的仪器设备、用具及其使用（强烈建议到实验室预习），完成课前预习思考题；了解实验需要测量的物理量，并根据要求提前准备实验记录表格（第一循环实验已由教师提供模板，可以打印）。预习成绩低于10分（共20分）者不能做实验。
	    \item 实验记录：认真、客观记录实验条件、实验过程中的现象以及数据。实验记录请用珠笔或者钢笔书写并签名（\textcolor{red}{\textbf{用铅笔记录的被认为无效}}）。\textcolor{red}{\textbf{保持原始记录，包括写错删除部分，如因误记需要修改记录，必须按规范修改。}}（不得输入电脑打印，但可扫描手记后打印扫描件）；离开前请实验教师检查记录并签名。
	    \item 分析讨论：处理实验原始数据（学习仪器使用类型的实验除外），对数据的可靠性和合理性进行分析；按规范呈现数据和结果（图、表），包括数据、图表按顺序编号及其引用；分析物理现象（含回答实验思考题，写出问题思考过程，必要时按规范引用数据）；最后得出结论。
	\end{enumerate}
	\textbf{实验报告就是将预习报告、实验记录、和数据处理与分析合起来，加上本页封面。}
	\item 每次完成实验后的一周内交\textbf{实验报告}（特殊情况不能超过两周）。
	% \item 实验报告注意事项
	% 	\begin{enumerate}[label=\roman*.]
	% 		\item 
	% 	\end{enumerate}
\end{enumerate}


\clearpage
\tableofcontents
\clearpage

\setcounter{section}{0}
\section{D1-A \quad 电阻热噪声及玻尔兹曼常量测量 \ 预习报告}
	
\subsection{实验目的}
\begin{enumerate}
	\item 通过实验了解热噪声的特征，理解热噪声的本质及其统计特征；
	\item 测出玻尔兹曼常量；
	\item 学习 VISA 接口协议，通过 LabVIEW 设置锁相放大器参数和进行数据采集。
	
\end{enumerate}

\subsection{仪器用具}
\begin{table}[htbp]
	\centering
	\renewcommand\arraystretch{1.6}
	% \setlength{\tabcolsep}{10mm}
	\begin{tabular}{p{0.05\textwidth}|p{0.20\textwidth}|p{0.05\textwidth}|p{0.5\textwidth}}
	\hline
	编号& 仪器用具名称 & 数量 &  主要参数（型号，测量范围，测量精度等） \\
	\hline
	1	&	锁相放大器 	&1 	& OE1022（ 系列）\\

	2	&	示波器 	&	1 	& RIGOL DS2202A		 \\
	
	3	&	数字多用表 & 1 &	RIGOL DM3058E、 	\\
	
	4	&	测温仪	&1 & MS6514	\\
	
	5	&	PC 机	&	 & LabVIEW 环境及有 VISA 接口协议	\\

	6	&	可变温样品盒（带BNC 接头）&	1	&	6.6k$\Omega$、 13k$\Omega$、 41k$\Omega$、 66k$\Omega$、 98k$\Omega$、144k$\Omega$	选其一	\\
	\hline
\end{tabular}
\end{table}


\subsection{实验原理}

	% \textbf{实验方案先前已提交至seelight，故在此不作展示}

	\begin{enumerate}
		\item 约翰逊-奈奎斯特噪声
		
		\textbf{约翰逊-奈奎斯特噪声（Johnson–Nyquist noise，即我们熟知的热噪声）}是由于电导体内的电荷载流子（通常为电子）在平衡状态下的热扰动所产生的电子噪声。
		
		这种现象与施加电压无关。热噪声存在于所有电路中，在某些灵敏的电子设备（如无线电接收器）中，热噪声会淹没弱信号，成为限制电子测量仪器灵敏度的主要因素。热噪声与绝对温度成正比，因此一些灵敏的电子设备（如射电望远镜的接收器）被冷却至低温状态以改善其信噪比。
		
		热噪声的通用统计物理推导称为\textbf{涨落-耗散定理}（详见后），该定理使用广义的电阻抗或广义的电极化率来表征介质的性质。
		
		% \begin{figure}[htbp]
		% 	\centering
		% 	\includegraphics[width=0.5\linewidth]{images/Oscilloscope-setup-johnson-noise}
		% 	\caption{约翰逊于1927年的实验表明，如果温度为$T$的电阻$R$的热噪声被限制在带宽$\Delta f$内，则其均方根电压($V_{\text{rms}}$)一般为$\sqrt{4 k_\text{B} T R \Delta f}$，其中$k_\text{B}$为玻尔兹曼常数。（来自\cite{a}）}
		% 	\label{fig:oscilloscope-setup-johnson-noise}
		% \end{figure}
		
		% \begin{figure}[htbp]
		% 	\centering
		% 	\includegraphics[width=0.75\linewidth]{images/Johnson-nyquist-noise-power-spectral-density-of-ideal-resistor}
		% 	\caption{约翰逊-奈奎斯特噪声的功率谱密度在每单位频率下接近常数，即$4k_\text{B}T R$，但由于量子效应，在高频（室温下为THz）时功率谱密度会衰减至零。该图的横轴使用对数刻度，每条竖线对应频率（以Hz为单位）的10次幂。（来自\cite{a}）}
		% 	\label{fig:johnson-nyquist-noise-power-spectral-density-of-ideal-resistor}
		% \end{figure}
		
		\begin{enumerate}
			\item 中等频率的理想电阻噪声
			
			约翰逊的实验发现，带宽为$\Delta f$的频带限制内，开氏温度$T$下电阻$R$的热噪声具有以下均方电压：
				$$
					\overline {V_n^2} = 4 k_\text{B} T R \, \Delta f
				$$
				其中$k_{\rm B}$是玻尔兹曼常数（$1.380649 \times 10^{-23}$ J/K）。

			该方程适用于非极端频率和温度下的理想电阻（即没有频率依赖性的纯电阻），但更精确的\textbf{通用形式}会考虑到\textbf{复阻抗}和量子效应。
			
			% \begin{figure}[htbp]
			% 	\centering
			% 	\includegraphics[width=0.4\linewidth]{images/Effect-of-bandwidth-on-selecting-noise}
			% 	\caption{尽管热噪声具有接近常数的功率谱密度$4 k_\text{B} T R$，但带宽为$\Delta f = f_\text{upper} - f_\text{lower}$的带通滤波器只通过高度为$4 k_\text{B} T R$、宽度为$\Delta f$的阴影部分。注意：实际的滤波器没有理想的“砖墙”截止频率，因此该区域的左右边缘并不是完全垂直的。（来自\cite{a}）}
			% 	\label{fig:effect-of-bandwidth-on-selecting-noise}
			% \end{figure}
			
			每赫兹带宽的均方电压为$4 k_\text{B} T R$，并且可以称为\textbf{功率谱密度}。在室温（约300 K）下，其平方根近似为$0.13 \sqrt{R}$，单位为$\mathrm{nV}/\sqrt{\mathrm{Hz}}$。例如，一个10 kΩ的电阻在室温下大约具有13 $\mathrm{nV}/\sqrt{\mathrm{Hz}}$。
			
			% \begin{figure}[htbp]
			% 	\centering
			% 	\includegraphics[width=0.3\linewidth]{images/JohnsonNoiseEquivalentCircuits}
			% 	\caption{这些电路是等效的：
			% 		\textbf{(A)} 带有内部热噪声的非零温度电阻；
			% 		\textbf{(B)} 其\textbf{戴维南等效}电路：一个无噪声电阻与噪声电压源串联；
			% 		\textbf{(C)} 其\textbf{诺顿等效}电路：一个无噪声电阻与噪声电流源并联。（来自\cite{a}）}
			% 	\label{fig:johnsonnoiseequivalentcircuits}
			% \end{figure}
			
			均方电压的平方根给出了在带宽$\Delta f$上的\textbf{均方根}（RMS）电压：
			\begin{equation*}
				V_\text{rms} = \sqrt{\overline {V_n^2}} = \sqrt{ 4 k_\text{B} T R \, \Delta f } \, .
			\end{equation*}
			
			带有热噪声的电阻可以用其\textbf{戴维南等效电路}表示，包含一个与高斯噪声电压源串联的无噪声电阻，其RMS电压如上所述。
			
			在室温下，3 kΩ的电阻在20 kHz（人类听力范围）带宽上产生的RMS噪声电压几乎为1微伏，$R \, \Delta f$为60 Ω·Hz时，RMS噪声电压几乎为1纳伏。
			
			带有热噪声的电阻还可以\textbf{转换为诺顿等效电路}，其等效电路为一个与高斯噪声电流源并联的无噪声电阻，其均方根电流为：
			\begin{equation*}
				I_\text{rms} = {V_\text{rms} \over R} = \sqrt {{4 k_\text{B} T \Delta f } \over R} \, .
			\end{equation*}
			
			\item 奈奎斯特对理想电阻噪声的推导
			
			% \begin{figure}[htbp]
			% 	\centering
			% 	\includegraphics[width=0.4\linewidth]{images/Nyquist-transmission-line-derivation-of-johnson-noise-with-voltage-sources}
				
			% 	\caption{奈奎斯特于1928年提出的思想实验的示意图。他使用两个有噪声的电阻（每个电阻由无噪声电阻和噪声电压源串联表示）通过一条长度为$l$的无损传输线连接。每个电阻的噪声信号以速度$v$沿传输线传播。由于所有阻抗相同，因此两个信号均被对方的电阻吸收而不是反射。奈奎斯特接着设想将传输线两端短路，从而将在线上运行的能量困住。由于阻抗不匹配，所有在途能量完全反射，此时可以用正弦驻波的叠加来表示这部分能量。在频带$\Delta f$内，有$2l \, \Delta f / v$个振荡模式。每个模式平均提供$k_\text{B} T$的能量，其中一半为电能（$k_\text{B} T / 2$），另一半为磁能，因此该频带内的总能量平均为$k_\text{B} T \cdot 2l \, \Delta f / v$。每个电阻贡献了$k_\text{B} T \cdot l \, \Delta f / v$的能量。由于短路前没有反射，传输线上总的在途能量也等于从两个电阻转移到传输线上的能量之和。除以传播时间间隔$l / v$，每个电阻平均在带宽$\Delta f$上转移的总功率为$k_\text{B} T \, \Delta f$。（来自\cite{a}）}
			% 	\label{fig:nyquist-transmission-line-derivation-of-johnson-noise-with-voltage-sources}
			% \end{figure}
			
			奈奎斯特在1928年的论文《导体中电荷的热激发》中，使用了玻尔兹曼和麦克斯韦的能量均分定律中关于势能和谐振子的概念，来解释约翰逊的实验结果。奈奎斯特的\textbf{思想实验}将两个相等电阻（$R_1 = R_2$）之间长无损传输线上的每个驻波模式的能量进行了求和。
			
			$R_1$在带宽$\Delta f$上传输到$R_2$的平均功率为：
			\begin{equation*}
				\overline{P_1} = k_\text{B} T \, \Delta f \, .
			\end{equation*}
			
			$V_1$（仅$R_1$的热噪声电压）通过组合电阻的电流为：
			\begin{equation*}
				I_1 = \frac{V_1}{R_1 + R_2} = \frac{V_1}{2R_1} \, ,
			\end{equation*}
			
			因此，$R_1$传输到$R_2$的功率为该电流的平方乘以$R_2$，简化为：
			\begin{equation*}
				P_1 = I_1^2 R_2 = I_1^2 R_1 = \left( \frac{V_1}{2R_1} \right)^2 R_1 = \frac{V_1^2}{4R_1} \, .
			\end{equation*}
			
			将$P_1$等于之前的平均功率表达式$\overline{P_1}$，可以求出在该带宽上的$V_1^2$的平均值：
			\begin{equation*}
				\overline{V_1^2} = 4 k_\text{B} T R_1 \, \Delta f \, .
			\end{equation*}
			
			\item 涨落耗散定理与电阻器中的热噪声
			
			涨落-耗散定理可以通过多种方式来表述，其中一种特别有用的形式如下：设$x(t)$是一个具有哈密顿量$H_0(x)$的动力系统的\textbf{可观测量}，该系统受到热涨落的影响。可观测量$x(t)$将在其平均值$\langle x\rangle_0$附近波动，这些波动可以通过\textbf{功率谱}$S_x(\omega) = \langle \hat{x}(\omega)\hat{x}^*(\omega) \rangle$来描述。假设我们可以施加一个时间变化的、空间上恒定的场$f(t)$，该场使哈密顿量变为$H(x) = H_0(x) - f(t)x$。可观测量$x(t)$对时间依赖场$f(t)$的响应，可以通过系统的\textbf{磁化率}或\textbf{线性响应函数}$\chi(t)$的一级近似来描述：
			\begin{equation*}
				\langle x(t) \rangle = \langle x \rangle_0 + \int_{-\infty}^{t} f(\tau) \chi(t-\tau) \, d\tau,
			\end{equation*}
			其中，摄动在$\tau = -\infty$时被逐渐（非常缓慢地）引入。
			
			涨落-耗散定理将$x$的双边功率谱（即正频率和负频率）与磁化率$\chi(t)$的傅里叶变换$\hat{\chi}(\omega)$的虚部联系起来：
			\begin{equation*}
				S_x(\omega) = -\frac{2 k_\mathrm{B} T}{\omega} \operatorname{Im}\hat{\chi}(\omega).
			\end{equation*}
			该等式成立的前提是傅里叶变换的约定为$f(\omega) = \int_{-\infty}^{\infty} f(t) e^{-i\omega t} \, dt$。左侧描述了$x$的\textbf{涨落}，右侧则与系统在受到振荡场$f(t) = F \sin(\omega t + \phi)$激励时\textbf{耗散}的能量密切相关。涨落的频谱揭示了线性响应，因为过去的涨落通过线性响应引起未来的涨落。
			
			约翰逊-奈奎斯特噪声，均方电压取决于电阻$R$、$k_{\rm B}T$以及测量电压的带宽$\Delta\nu$：
			\begin{equation*}
				\langle V^2 \rangle \approx 4Rk_{\rm B}T\,\Delta\nu \, .
			\end{equation*}
			
			这一观察可以通过\textbf{涨落-耗散定理}来理解。以一个由电阻$R$和小电容$C$组成的简单电路为例。根据\textbf{基尔霍夫电路定律}，电压满足以下关系：
			\begin{equation*}
				V = -R\frac{dQ}{dt} + \frac{Q}{C} \, .
			\end{equation*}
			
			因此，该电路的\textbf{响应函数}为
			\begin{equation*}
				\chi(\omega) \equiv \frac{Q(\omega)}{V(\omega)} = \frac{1}{\frac{1}{C} - i\omega R} \, .
			\end{equation*}
			
			在低频极限$\omega \ll (RC)^{-1}$下，响应函数的虚部近似为
			\begin{equation*}
				\text{Im}\left[\chi(\omega)\right] \approx \omega RC^2 \, .
			\end{equation*}
			
			通过涨落-耗散定理，可以将其与电压的\textbf{功率谱密度}函数$S_V(\omega)$联系起来
			\begin{equation*}
				S_V(\omega) = \frac{S_Q(\omega)}{C^2} \approx \frac{2k_{\rm B}T}{C^2\omega} \text{Im}\left[\chi(\omega)\right] = 2Rk_{\rm B}T \, .
			\end{equation*}
			
			在带宽$\Delta f = \Delta\omega / (2\pi)$内观察到的约翰逊-奈奎斯特电压噪声$\langle V^2 \rangle$，中心频率为$\omega = \pm \omega_0$。因此
			\begin{equation*}
				\langle V^2 \rangle \approx S_V(\omega) \times 2\Delta f \approx 4Rk_{\rm B}T\Delta f \, .
			\end{equation*}
			
		\end{enumerate}
		
		\item 噪声的测量
		
		实验利用锁相放大器测量噪声。时域上做等时间间隔采样，多次测量计算信号的\textbf{方均根值}。这句话的意思是：设备在等间隔的时间点上读取一个信号 $v(t_i) = s(t_i) + n(t_i)$，$s(t_i)$ 为 $t_i$ 时刻频率为 $f$ 的信号，$n(t_i)$ 为 $t_i$ 时刻的噪声，经与信号同频率 $f$ 的参考信号相敏检波后，输出 X 分量$X_f(t_i) = s_x(t_i) + n_f(t_i)$，是包含“同频”噪声的信号。定义噪声强度 X 分量测量值（这个式子显然是信号的方均根的表达式，同时它表达了噪声电压有效值）：
		\[
		v_{Nx}(f) = \sqrt{\frac{\sum_{i=1}^{M} \left(X_f(t_i) - \bar{X_f}\right)^2}{M}}
		\]
		其中，$\bar{X}_j = \frac{1}{M} \sum_{j=1}^{M} X_j$为输入信号的平均值，$M$ 为用于计算输入噪声值的样本数。
		
		实际测量有各种各样的噪声，考虑按以下方式合成：
		\[
		v_{Nx}(f) = \sqrt{\frac{\sum_{i=1}^{M} \left( n_{f_x}(t_i) \right)^2}{M}} = \sqrt{\frac{\sum_{i=1}^{M} \left( \sum_{k=1}^{m} n_{f_xk}(t_i) \right)^2}{M}} 
		\]
		
		分离不同种类的噪声是困难的，除非在输入端做物理隔离。所以我们考虑的是非相干的噪声，按以下方式做合成：
		\[
		v_{Nx}(f) = \sqrt{\frac{\sum_{k=1}^{m} \sum_{i=1}^{M} \left( n_{f_xk}(t_i) \right)^2}{M}}
		\]
		上式给出了不同类型非相干噪声的叠加原理。
		
		对随机噪声，X 与 Y 独立等价，即输入噪声功率为 X 或 Y 测量功率的 2 倍：
		\[
		v_N^2 = \frac{1}{M} \sum_{j=1}^{M} \left( n_{xT}^2(j) + n_{yT}^2(j) \right) = \frac{2}{M} \sum_{j=1}^{M} n_{xT}^2(j)
		\]
		
		\item 仪器本底扣除及输入带宽修正
		
		测量电阻 $R$ 的噪声 ($n_R$) 是待测热噪声 ($n_T$) 与本底噪声的叠加。根据随机噪声叠加原理式得：
		\[
		n_{Rx}^2(\omega) = n_{Tx}^2(\omega) + n_{Bx}^2(\omega)
		\]
		
		可扣除仪器的本底噪声，\textbf{前提是测量本底噪声时与测量噪声时的条件，以及所选的锁相放大器参数相同}。
		
		事实上，在后续进一步的资料查找阅读过程中，从\cite{g}中了解到如下信息，这进一步对电阻阻值和测温点的参数值的选取提供了参考：
		
		锁相放大器的本底噪声谱密度在 $5 \, \text{nV}/\sqrt{\text{Hz}}$ 左右， 相当于 $6 \, \text{k}\Omega$ 的电阻在液氮温度 (77 K) 下，或 $1.5 \, \text{k}\Omega$ 的电阻在室温 (300 K) 下的热噪声谱密度。可见，要降低热噪声测量的系统误差，一方面要求保证样本量足够大，另一方面电阻值选 $5 \sim 100 \, \text{k}\Omega$ 比较合适。
		
		进一步再考虑输入带宽修正：
		\[
		n_{Tx}^2 = n_{Rx}^2 \left[ 1 + (\omega R_{\text{in}} C_{\text{in}})^2 \right] - n_{Bx}^2
		\]
		其中，$R_{\text{in}}$ 和 $C_{\text{in}}$ 分别为锁相放大器的输入端的等效电阻和等效电容。
		
		带宽修正项 \(1 + (\omega R_{\text{in}}C_{\text{in}})^2\) 反映了输入阻抗和电容对信号频率响应的影响，这个修正项可以通过输入电路的频率响应特性推导出来，特别是输入阻抗 \( R_{\text{in}} \) 和电容 \( C_{\text{in}} \) 所形成的低通滤波效应。锁相放大器的输入端通常有输入电阻和输入电容，这两者并联形成了一个简单的RC电路，对输入信号产生低通滤波作用。假设输入信号电压为 \( V_{\text{in}} \)，在通过该RC电路时会产生频率依赖的输出电压 \( V_{\text{out}} \)。这一过程可以用传递函数来描述为：
		\[
		H(\omega) = \frac{V_{\text{out}}}{V_{\text{in}}} = \frac{1}{1 + j\omega R_{\text{in}} C_{\text{in}}}
		\]
		其中 \( \omega = 2\pi f \) 为角频率，\( j \) 为虚数单位。
		
		为了得到输入信号经过该RC电路后的衰减效果，我们需要计算传递函数的幅值平方：
		\[
		|H(\omega)|^2 = \left| \frac{1}{1 + j\omega R_{\text{in}} C_{\text{in}}} \right|^2 = \frac{1}{1 + (\omega R_{\text{in}} C_{\text{in}})^2}
		\]
		这个结果表示信号经过RC电路后，随频率增高而衰减。
		
		测量电阻热噪声时，电阻产生的噪声信号通过该RC电路，因此噪声的频谱功率密度也会随着频率而衰减。由于热噪声是白噪声，在不同频率下均匀分布，我们需要对测量到的噪声功率进行修正。假设 \( n_{\text{Rx}}^2 \) 是未经修正的噪声功率，那么经过RC电路后的噪声功率会变为：
		\[
		n_{\text{Tx}}^2 = n_{\text{Rx}}^2 \cdot |H(\omega)|^2 = n_{\text{Rx}}^2 \cdot \frac{1}{1 + (\omega R_{\text{in}} C_{\text{in}})^2}
		\]
		
		由于实验的目标是测量未经滤波器影响的噪声信号，因此我们需要将传递函数中的频率衰减效应修正回来。通过乘以传递函数的倒数，可以得到修正后的噪声功率：
		\[
		n_{\text{Tx}}^2 = n_{\text{Rx}}^2 \cdot \left[ 1 + (\omega R_{\text{in}} C_{\text{in}})^2 \right]
		\]
		这样我们就补偿了输入电路对噪声信号的衰减效应，得到接近理想条件下的噪声功率。
		
		此外，锁相放大器自身产生的本底噪声也需要从测量结果中扣除，假设本底噪声的功率为 \( n_{\text{Bx}}^2 \)，最终修正后的噪声功率为：
		\[
		n_{\text{Tx}}^2 = n_{\text{Rx}}^2 \left[ 1 + (\omega R_{\text{in}} C_{\text{in}})^2 \right] - n_{\text{Bx}}^2
		\]
		
		为什么本底噪声不必做输入带宽修正？本底噪声是锁相放大器内部的固有噪声，不经过外部输入电路，因此不会受到输入带宽的影响，也无需进行带宽修正。而外部噪声，如电阻的热噪声，会受到输入阻抗和滤波器的影响，因此需要进行带宽修正。
		
		此外，注意讲义上提到的非随机的环境噪声$n_{ex}$。
		
		实验讲义上提供了具体的数值：影响输入带宽的 $R_{\text{in}}$ 约为 $2/3R$，$C_{\text{in}}$ 约为电缆电容（室温下约为 96pF/m）、锁相放大器输入电容（25pF），与包括连接器和样品盒在内的电容并联后的值的 1/2（约 \textbf{68pF}）。
		
		\item $k_B$的计算
		
		首先考虑\textbf{等效噪声带宽}。实际滤波器带宽做不到理想的矩形，我们用谱线所包围的功率面积等于矩形面积来定义等效噪声带宽，等效噪声带宽 $\text{ENBW} = \Delta f_{\text{LPF}}/2$ 。将随机变化的噪声看作幅值随时间的阶跃输入信号，则窄带滤波器使输出信号延迟。定义输出信号达到阶跃的 99\% 视为反映输入噪声幅值，其延迟或等待时间也列于\cref{tab:ENBW}。$X$ 稳定需要的等待时间为 $t_X$。
		
		\begin{table}[h!]
			\centering
			\caption{单分量等效噪声带宽（ENBW）与陡降和时间常数$\tau$的对应关系}
			\label{tab:ENBW}
			\begin{tabular}{cccc}
				\hline
				\textbf{阶数(n)} & \textbf{Slope (dB/oct)} & \textbf{ENBW (1/$\tau$)} & \textbf{X 测量等待时间$t_X$}\\
				\hline
				1 & 6 & 1/4 & 4.6 \\
				2 & 12 & 1/8 & 6.6 \\
				3 & 18 & 3/32 & 8.4 \\
				4 & 24 & 5/64 & 10 \\
				\hline
			\end{tabular}
		\end{table}
		
		\textbf{注意ENBW需要通过计算得来，根据不同的陡降得到不同的系数（来自讲义的表），在乘以时间常数的倒数便得到带宽。}
		
		引入ENBW后，我们可以得到最终\textbf{玻尔兹曼常数的计算式}。考虑到，$v_N$ 测量值与仪器带宽有关，故用与仪器参数无关的噪声谱密度描述噪声特性应：
			$$
			S_T(f)=v_{Nx}/\sqrt{\text{ENBW}} = \sqrt{4k_BTR} \Rightarrow k_B = \dfrac{v_{Nx}^2}{4\times\text{ENBW}\times TR}
			$$
		
		作为参考，1.0k$\Omega$金属电阻在室温 (300K) 下噪声谱密度为：
		\[
		\sqrt{4 \times 1.38 \times 10^{-23} \text{J/K} \times 300\text{K} \times 1000\Omega} = 4.07\text{nV}/\sqrt{\text{Hz}}
		\]
		
		此外，基于$k_B = \dfrac{v_{Nx}^2}{4\times\text{ENBW}\times TR}$，我们可以考虑一种变温测量下的\textbf{线性拟合方案}，即用变化的温度与阻值（阻值也会随温度变化）乘积作为梯度测量下的功率谱变化的斜率以表征玻尔兹曼常数：
		$$
		\dfrac{v_{Nx}^2}{\text{ENBW}}=4k_B\times TR
		$$
		
		实际测量 $S_T(f)$ 时，锁相放大器采样次数不可能无限大，热噪声测量值的比值标准差与样本量的关系约为 $2/\sqrt{M}$。
		
		此外，OE1022 锁相放大器用式计算噪声电压谱密度（$N$ 值）的样本量为 1024 个。受限于低通滤波器，所测得的邻近两个 $N$ 值不一定是独立无关的，从而电压数据也未必服从正态分布。因此，\textbf{建议实验时直接记录 X 或 Y，再计算噪声强度和噪声谱密度更为清晰可靠。}
		
	\end{enumerate}

	\vspace{1cm}
	\textbf{\huge 更具体的实验方案先前已上传至seelight}






\clearpage
\subsection{实验前思考题}

% 思考题1
\begin{question}
	如何与同班组的同学分配频率点？
\end{question}


	在实验中合理分配频率点有助于提高效率，避免资源浪费。首先需确定测量的频率范围和步进，通常依据实验要求和设备性能选择，例如1Hz至100kHz。分配方式可以有线性划分和对数划分：

	线性划分：将频率点在1Hz到100kHz之间均匀分布。例如，测量20个频率点，若每组4人，每人负责5个频率点。这种方式简单、均匀，确保每人覆盖从低频到高频的频率段。

	对数划分：适用于跨度较大的频率范围。在对数坐标上均匀分配，如1Hz、10Hz、100Hz等倍数变化的频率点，这样可以更好地覆盖高频段变化，同时保持低频的合理间隔。

	为了确保数据一致性，建议在一些频率点上进行交叉验证。例如，同学A和B在10Hz和1kHz重复测量，C和D在50Hz和5kHz交叉测量，以提高实验数据的可靠性。





\begin{question}
	在测量所选参量前或后，是否需要对锁相放大器进行校正？如何校正？
\end{question}


	在电阻热噪声和玻尔兹曼常量测量实验中，校正锁相放大器是确保测量准确性的关键。主要校正步骤如下：

	零点校正：在没有输入信号时，锁相放大器可能显示输出电压，这是由于零点漂移导致的。为校正零点，应短接或断开输入端，调整放大器的偏移设置，使输出保持为零，从而减少零点漂移引入的误差。

	本底噪声校正：锁相放大器存在热噪声、温漂噪声等本底噪声。测量外部信号前，先测量设备的本底噪声平方值 $n_{Bx}^2$。测量总噪声平方值 $n{Rx}^2$ 后，使用公式 $n_{Tx}^2 = n_{Rx}^2 - n_{Bx}^2$扣除设备本底噪声，从而获得真实噪声信号。

	同时考虑输入带宽校正，有：
	\[
		n_{Tx}^2 = n_{Rx}^2[1 + (\omega R_{in} C_{in})^2] - n_{Bx}^2
	\]







\begin{question}
	如何辨别在某频率下测量到的噪声是电阻本身的热噪声还是带有环境噪声？
\end{question}

	\begin{enumerate}
		\item 频率特性分析：电阻的热噪声是白噪声，功率谱密度平坦；环境噪声（如市电干扰）具有特定频率成分，常见于50Hz及其倍频。通过观察频谱中的峰值频率，可识别出环境噪声成分。
		\item 本底噪声校正：锁相放大器本身有本底噪声（如设备热噪声）。为获得电阻的真实热噪声，需在断开信号输入时测量本底噪声，并在总噪声中扣除该值，从而减少设备对实验结果的干扰。
		\item 带宽调整：热噪声强度随带宽线性增加，若噪声信号随带宽线性变化，则主要为热噪声；若变化不线性或波动较大，可能是环境噪声的影响。
		\item 屏蔽干扰：使用法拉第笼等屏蔽措施有效隔离外界电磁干扰，可减少环境噪声影响，使测量结果更接近真实的热噪声。
		\item 实验时间选择：避免环境噪声影响，可选择在夜间或设备较少的环境中进行测量，减少干扰确保测量的噪声主要来自电阻的热噪声。
	\end{enumerate}
	

	

	

	

	





\clearpage
\begin{table}
	\renewcommand\arraystretch{1.7}
	\centering
	\begin{tabularx}{\textwidth}{|X|X|X|X|}
	\hline
	专业：& 物理学 &年级：& 2022级 \\
	\hline
	姓名：& 戴鹏辉 & 学号：& 22344016 \\
	\hline
	室温：& 26℃ & 实验地点： & A101 \\
	\hline
	学生签名：& \includegraphics[width=0.9in]{name-TaLEs.jpg} & 评分： &\\
	\hline
	实验时间：& 2024/10/25 & 教师签名：&\\
	\hline
	\end{tabularx}
\end{table}

\section{D1 \quad 电阻热噪声及玻尔兹曼常量测量 \quad\heiti 实验记录}
\subsection{实验内容、步骤和数据记录}

	\subsubsection*{实验1.6 \quad 滤波器带宽对输出信号响应的影响}
	
	\begin{figure}[H]
		\centering
		% 第一行左图
		\begin{minipage}[b]{0.40\textwidth}
			\centering
			\includegraphics[width=\textwidth]{30us 24.jpg} % 替换为你的图片路径
			\caption*{$\tau = 30 \mu s \ Slope = 24 \mathrm{dB/oct}$}
			\label{fig:D1-6-1-1}
		\end{minipage}
		\hspace{0.05\textwidth} % 两列之间的水平间距
		% 第一行右图
		\begin{minipage}[b]{0.40\textwidth}
			\centering
			\includegraphics[width=\textwidth]{100us 24.jpg} % 替换为你的图片路径
			\caption*{$\tau = 100 \mu s \ Slope = 24 \mathrm{dB/oct}$}
			\label{fig:D1-6-1-2}
		\end{minipage}
		
		\vspace{0.01\textwidth} % 两行之间的垂直间距
	
		% 第二行左图
		\begin{minipage}[b]{0.40\textwidth}
			\centering
			\includegraphics[width=\textwidth]{300us 24.jpg} % 替换为你的图片路径
			\caption*{$\tau = 300 \mu s \ Slope = 24 \mathrm{dB/oct}$}
			\label{fig:D1-6-1-3}
		\end{minipage}
		\hspace{0.05\textwidth} % 两列之间的水平间距
		% 第二行右图
		\begin{minipage}[b]{0.40\textwidth}
			\centering
			\includegraphics[width=\textwidth]{1ms 24.jpg} % 替换为你的图片路径
			\caption*{$\tau = 1 ms \ Slope = 24 \mathrm{dB/oct}$}
			\label{fig:D1-6-1-4}
		\end{minipage}
		\caption{相同陡降、不同时间常数的信号响应}
		\label{fig:D1-6-1}
	\end{figure}


	
	由\cref{fig:D1-6-1}，在保持陡降不变的情况下，随着时间常数从30µs增加到1ms，信号变化的幅度会逐渐减小。
	
	这是因为较大的时间常数使低通滤波器的通带变窄，从而抑制了高频成分，导致输出信号对快速变化的输入信号响应变得迟缓。
	
	
	

	
	\begin{figure}[H]
		\centering
		% 第一行左图
		\begin{minipage}[b]{0.40\textwidth}
			\centering
			\includegraphics[width=\textwidth]{100us 6.jpg} % 替换为你的图片路径
			\caption*{$\tau = 100 \mu s \ Slope = 6 \mathrm{dB/oct}$}
			\label{fig:D1-6-1-1}
		\end{minipage}
		\hspace{0.05\textwidth} % 两列之间的水平间距
		% 第一行右图
		\begin{minipage}[b]{0.40\textwidth}
			\centering
			\includegraphics[width=\textwidth]{100us 12.jpg} % 替换为你的图片路径
			\caption*{$\tau = 100 \mu s \ Slope = 12 \mathrm{dB/oct}$}
			\label{fig:D1-6-1-2}
		\end{minipage}
		
		\vspace{0.01\textwidth} % 两行之间的垂直间距
	
		% 第二行左图
		\begin{minipage}[b]{0.40\textwidth}
			\centering
			\includegraphics[width=\textwidth]{100us 18.jpg} % 替换为你的图片路径
			\caption*{$\tau = 100 \mu s \ Slope = 18 \mathrm{dB/oct}$}
			\label{fig:D1-6-1-3}
		\end{minipage}
		\hspace{0.05\textwidth} % 两列之间的水平间距
		% 第二行右图
		\begin{minipage}[b]{0.40\textwidth}
			\centering
			\includegraphics[width=\textwidth]{100us 24.jpg} % 替换为你的图片路径
			\caption*{$\tau = 1 \mu s \ Slope = 24 \mathrm{dB/oct}$}
			\label{fig:D1-6-1-4}
		\end{minipage}
		\caption{相同时间常数、不同陡降的信号响应}
		\label{fig:D1-6-2}
	\end{figure}

	由\cref{fig:D1-6-2}，在保持时间常数不变的情况下，当陡降从6dB/oct增加到24dB/oct时，信号变化的幅度逐渐减小。因为陡降的增加会对高频噪声的抑制更加显著。
	
	
	对于不同变化特性的待测信号，应合理选择时间常数和陡降值，以达到最佳的信号捕捉和处理效果。增大时间常数有助于提高信号的稳定性，但并非越大越好，合适的时间常数能保证系统具备足够的响应能力，准确捕捉$R(t)$的变化。同时，适当的陡降设置能够有效抑制噪声，并确保信号变化过程的准确显示，避免过高的陡降导致系统对快速信号变化的敏感度降低，影响测量精度。因此，在时间常数和陡降之间找到平衡至关重要。











	\clearpage
	\subsubsection*{电阻热噪声及玻尔兹曼常量测量}

		\begin{enumerate}
			\item 首先，我们选择了标记为 144kΩ的电阻进行电阻热噪声的测量。已知样品盒采用3个1/8W的阻值为R的色环电阻，以“三角”对称连接方式解决因锁相放大器输入电阻所带来的输入带宽变窄的问题，其等效电阻为2/3R。
			\item 线路连接上，为尽量抑制环境噪声的影响，测量选用差分（A-B）输入，电缆BNC接头分别连接在锁相放大器的A 和B端。
			\item 仪器参数方面，满量程灵敏度应选为1µV；另一方面，热噪声幅值小，动态储备应选“Low Noise”。
			
			由于测量装置不可能完全屏蔽环境干扰，选择太大的带宽容易引入近频的环境噪声，因此建议Slope > 12dB/oct。

			为保证所采集的每个数据是独立的并且不失真，应该选择合适的采样的时间间隔。在我们的具体实验中，选择时间常数$\tau$ 为30ms、Slope 为 18 dB/oct ，ENBW 为 3/32 Hz，X 值的测量等待时间或时间间隔约为8.4$\tau$ ，X，Y 值同时采样 $10^3$ 个需要时间约为4.5 分钟。
			\item 随后进行本底测量，用50Ω短接，以进行本地校正。
		\end{enumerate}


		我们参照实验方案，分别在常温（298K）、和低温（74.3K）下，测量了电阻的热噪声，并由此计算了玻尔兹曼常数。得到的数据如\cref{tbl:D1-A-1}：


		实验软件的界面截图如\cref{fig:PrintScreen}：





		\begin{table}[htbp]
			\centering
			\begin{tblr}{
			  cells = {c},
			  vline{1-2,4,7} = {-}{},
			  hline{1-2,10} = {-}{},
			}
			频率/Hz     & 时间常数/ms & 陡降/(dB/oct) & 本底数据存储文件名 & 热噪声数据储文件名 & 备注  \\
			6666  & 10   & 18 & DataB1    & Data1    & 常温1 \\
			1145  & 10   & 18 & DataB2    & Data2    & 常温2 \\
			11451 & 10   & 18 & DataB3    & Data3    & 常温3 \\
			20486 & 10   & 18 & DataB4    & Data4    & 常温4 \\
			666   & 10   & 18 & DataB5    & Data5    & 低温1 \\
			6666  & 10   & 18 & DataB6    & Data6    & 低温2 \\
			66666 & 10   & 18 & DataB7    & Data7    & 低温3 \\
			36666 & 10   & 18 & DataB8    & Data8    & 低温4 
			\end{tblr}
			\caption{本底和热噪声测量结果记录表}
			\label{tbl:D1-A-1}
		\end{table}

		在我们的实验过程中，我们选择保持时间常数和陡降不变，即$\tau = 10 \mathrm{ms}$，$Slope = 18 \mathrm{dB/oct}$，采样时间间隔取$\Delta t = 10 \tau = 100 \mathrm{ms}$


		具体的分析讨论与数据处理过程见后文以及附件中的 jupyternotebook。




		\begin{figure}[htbp]
			\centering
			\begin{minipage}{0.45\textwidth}
				\subfloat[]{\includegraphics[width=\linewidth]{PrintScreen1.png}\label{fig:PrintScreen1}}
			\end{minipage}
			\hfill
			\begin{minipage}{0.45\textwidth}
				\subfloat[]{\includegraphics[width=\linewidth]{PrintScreen2.png}\label{fig:PrintScreen2}}
			\end{minipage}
			
			
			\caption{实验软件界面截图}
			\label{fig:PrintScreen}
		\end{figure}


















% \clearpage
\subsection{实验过程中遇到的问题记录}

\begin{enumerate}
	\item 注意应当先完成常温条件下的测试，再完成低温条件下的测试，因为电阻的升温降温需要较长时间。
	\item 注意每一次改变仪器参数的设置，都要重新测量本底噪声。
	
\end{enumerate}


\clearpage
\subsection{原始数据记录}

	\begin{figure}[H]
		\centering
		\subfloat[原始数据1]{\includegraphics[width=0.35\textwidth]{APL1-1-OriData1.jpg}\label{fig:data1}}
		\quad
		\subfloat[原始数据2]{\rotatebox{90}{\includegraphics[width=0.50\textwidth]{APL1-1-OriData2.jpg}}\label{fig:data2}}
		\quad
		\subfloat[原始数据3]{\includegraphics[width=0.35\textwidth]{APL1-1-OriData3.jpg}\label{fig:data3}}
		\quad

		\caption{原始数据}
		\label{fig:data}
	\end{figure}
	

\clearpage
\begin{table}
	\renewcommand\arraystretch{1.7}
	\begin{tabularx}{\textwidth}{|X|X|X|X|}
	\hline
	专业：& 物理学 &年级：& 2022级\\
	\hline
	姓名： & 戴鹏辉 & 学号：& 22344016\\
	\hline
    日期：& 2024/10/25 & 评分： &\\
	\hline
	\end{tabularx}
\end{table}

\section{D1 \quad 电阻热噪声及玻尔兹曼常量测量 \quad\heiti 分析与讨论}

\subsection{实验数据分析}

	\subsection{数据的质量评定}

		首先我们对测量得到的原始数据进行质量评定，以保证所采集的每个数据是独立的并且不失真，在保证数据优质的情况下，再往下进行玻尔兹曼常数的计算。

		数据分析如下，文件名对应的实验条件见\cref{tbl:D1-A-1}：
		% \begin{figure}[htbp]
		% 	\centering
		% 	% 第一行图片
		% 	\begin{subfigure}[b]{0.23\textwidth}
		% 		\includegraphics[width=\textwidth]{D1-A-3-Data1.png}
		% 		\caption{图片1}
		% 		\label{fig:image1}
		% 	\end{subfigure}
		% 	\hfill
		% 	\begin{subfigure}[b]{0.23\textwidth}
		% 		\includegraphics[width=\textwidth]{D1-A-3-Data2.png}
		% 		\caption{图片2}
		% 		\label{fig:image2}
		% 	\end{subfigure}
		% 	\hfill
		% 	\begin{subfigure}[b]{0.23\textwidth}
		% 		\includegraphics[width=\textwidth]{D1-A-3-Data3.png}
		% 		\caption{图片3}
		% 		\label{fig:image3}
		% 	\end{subfigure}
		% 	\hfill
		% 	\begin{subfigure}[b]{0.23\textwidth}
		% 		\includegraphics[width=\textwidth]{D1-A-3-Data4.png}
		% 		\caption{图片4}
		% 		\label{fig:image4}
		% 	\end{subfigure}
			
		% 	% 第二行图片
		% 	\begin{subfigure}[b]{0.23\textwidth}
		% 		\includegraphics[width=\textwidth]{D1-A-3-Data5.png}
		% 		\caption{图片5}
		% 		\label{fig:image5}
		% 	\end{subfigure}
		% 	\hfill
		% 	\begin{subfigure}[b]{0.23\textwidth}
		% 		\includegraphics[width=\textwidth]{D1-A-3-Data6.png}
		% 		\caption{图片6}
		% 		\label{fig:image6}
		% 	\end{subfigure}
		% 	\hfill
		% 	\begin{subfigure}[b]{0.23\textwidth}
		% 		\includegraphics[width=\textwidth]{D1-A-3-Data7.png}
		% 		\caption{图片7}
		% 		\label{fig:image7}
		% 	\end{subfigure}
		% 	\hfill
		% 	\begin{subfigure}[b]{0.23\textwidth}
		% 		\includegraphics[width=\textwidth]{D1-A-3-Data8.png}
		% 		\caption{图片8}
		% 		\label{fig:image8}
		% 	\end{subfigure}
			
		% 	% 第三行图片
		% 	\begin{subfigure}[b]{0.23\textwidth}
		% 		\includegraphics[width=\textwidth]{D1-A-3-DataB1.png}
		% 		\caption{图片9}
		% 		\label{fig:image9}
		% 	\end{subfigure}
		% 	\hfill
		% 	\begin{subfigure}[b]{0.23\textwidth}
		% 		\includegraphics[width=\textwidth]{D1-A-3-DataB2.png}
		% 		\caption{图片10}
		% 		\label{fig:image10}
		% 	\end{subfigure}
		% 	\hfill
		% 	\begin{subfigure}[b]{0.23\textwidth}
		% 		\includegraphics[width=\textwidth]{D1-A-3-DataB3.png}
		% 		\caption{图片11}
		% 		\label{fig:image11}
		% 	\end{subfigure}
		% 	\hfill
		% 	\begin{subfigure}[b]{0.23\textwidth}
		% 		\includegraphics[width=\textwidth]{D1-A-3-DataB4.png}
		% 		\caption{图片12}
		% 		\label{fig:image12}
		% 	\end{subfigure}
			
		% 	% 第四行图片
		% 	\begin{subfigure}[b]{0.23\textwidth}
		% 		\includegraphics[width=\textwidth]{D1-A-3-DataB5.png}
		% 		\caption{图片13}
		% 		\label{fig:image13}
		% 	\end{subfigure}
		% 	\hfill
		% 	\begin{subfigure}[b]{0.23\textwidth}
		% 		\includegraphics[width=\textwidth]{D1-A-3-DataB6.png}
		% 		\caption{图片14}
		% 		\label{fig:image14}
		% 	\end{subfigure}
		% 	\hfill
		% 	\begin{subfigure}[b]{0.23\textwidth}
		% 		\includegraphics[width=\textwidth]{D1-A-3-DataB7.png}
		% 		\caption{图片15}
		% 		\label{fig:image15}
		% 	\end{subfigure}
		% 	\hfill
		% 	\begin{subfigure}[b]{0.23\textwidth}
		% 		\includegraphics[width=\textwidth]{D1-A-3-DataB8.png}
		% 		\caption{图片16}
		% 		\label{fig:image16}
		% 	\end{subfigure}
			
		% 	\caption{16张图片}
		% 	\label{fig:images}
		% \end{figure}


		% \begin{figure}[htbp]
		% 	\centering
		% 	% 第一行图片
		% 	\subfloat[图片1]{\includegraphics[width=0.23\textwidth]{D1-A-3-Data1.png}\label{fig:image1}}
		% 	\hfill
		% 	\subfloat[图片2]{\includegraphics[width=0.23\textwidth]{D1-A-3-Data2.png}\label{fig:image2}}
		% 	\hfill
		% 	\subfloat[图片3]{\includegraphics[width=0.23\textwidth]{D1-A-3-Data3.png}\label{fig:image3}}
		% 	\hfill
		% 	\subfloat[图片4]{\includegraphics[width=0.23\textwidth]{D1-A-3-Data4.png}\label{fig:image4}}
			
		% 	% 第二行图片
		% 	\subfloat[图片5]{\includegraphics[width=0.23\textwidth]{D1-A-3-Data5.png}\label{fig:image5}}
		% 	\hfill
		% 	\subfloat[图片6]{\includegraphics[width=0.23\textwidth]{D1-A-3-Data6.png}\label{fig:image6}}
		% 	\hfill
		% 	\subfloat[图片7]{\includegraphics[width=0.23\textwidth]{D1-A-3-Data7.png}\label{fig:image7}}
		% 	\hfill
		% 	\subfloat[图片8]{\includegraphics[width=0.23\textwidth]{D1-A-3-Data8.png}\label{fig:image8}}
			
		% 	% 第三行图片
		% 	\subfloat[图片9]{\includegraphics[width=0.23\textwidth]{D1-A-3-DataB1.png}\label{fig:image9}}
		% 	\hfill
		% 	\subfloat[图片10]{\includegraphics[width=0.23\textwidth]{D1-A-3-DataB2.png}\label{fig:image10}}
		% 	\hfill
		% 	\subfloat[图片11]{\includegraphics[width=0.23\textwidth]{D1-A-3-DataB3.png}\label{fig:image11}}
		% 	\hfill
		% 	\subfloat[图片12]{\includegraphics[width=0.23\textwidth]{D1-A-3-DataB4.png}\label{fig:image12}}
			
		% 	% 第四行图片
		% 	\subfloat[图片13]{\includegraphics[width=0.23\textwidth]{D1-A-3-DataB5.png}\label{fig:image13}}
		% 	\hfill
		% 	\subfloat[图片14]{\includegraphics[width=0.23\textwidth]{D1-A-3-DataB6.png}\label{fig:image14}}
		% 	\hfill
		% 	\subfloat[图片15]{\includegraphics[width=0.23\textwidth]{D1-A-3-DataB7.png}\label{fig:image15}}
		% 	\hfill
		% 	\subfloat[图片16]{\includegraphics[width=0.23\textwidth]{D1-A-3-DataB8.png}\label{fig:image16}}
			
		% 	\caption{16张图片}
		% 	\label{fig:images}
		% \end{figure}

		% \begin{figure}[htbp]
		% 	\centering
		% 	% 第一行图片
		% 	\subfloat[图片1]{\includegraphics[width=0.45\textwidth]{D1-A-3-Data1.png}\label{fig:image1}}
		% 	\hfill
		% 	\subfloat[图片2]{\includegraphics[width=0.45\textwidth]{D1-A-3-Data2.png}\label{fig:image2}}
			
		% 	% 第二行图片
		% 	\subfloat[图片3]{\includegraphics[width=0.45\textwidth]{D1-A-3-Data3.png}\label{fig:image3}}
		% 	\hfill
		% 	\subfloat[图片4]{\includegraphics[width=0.45\textwidth]{D1-A-3-Data4.png}\label{fig:image4}}
			
		% 	% 第三行图片
		% 	\subfloat[图片5]{\includegraphics[width=0.45\textwidth]{D1-A-3-Data5.png}\label{fig:image5}}
		% 	\hfill
		% 	\subfloat[图片6]{\includegraphics[width=0.45\textwidth]{D1-A-3-Data6.png}\label{fig:image6}}
			
		% 	% 第四行图片
		% 	\subfloat[图片7]{\includegraphics[width=0.45\textwidth]{D1-A-3-Data7.png}\label{fig:image7}}
		% 	\hfill
		% 	\subfloat[图片8]{\includegraphics[width=0.45\textwidth]{D1-A-3-Data8.png}\label{fig:image8}}
			
		% 	% 第五行图片
		% 	\subfloat[图片9]{\includegraphics[width=0.45\textwidth]{D1-A-3-DataB1.png}\label{fig:image9}}
		% 	\hfill
		% 	\subfloat[图片10]{\includegraphics[width=0.45\textwidth]{D1-A-3-DataB2.png}\label{fig:image10}}
			
		% 	% 第六行图片
		% 	\subfloat[图片11]{\includegraphics[width=0.45\textwidth]{D1-A-3-DataB3.png}\label{fig:image11}}
		% 	\hfill
		% 	\subfloat[图片12]{\includegraphics[width=0.45\textwidth]{D1-A-3-DataB4.png}\label{fig:image12}}
			
		% 	% 第七行图片
		% 	\subfloat[图片13]{\includegraphics[width=0.45\textwidth]{D1-A-3-DataB5.png}\label{fig:image13}}
		% 	\hfill
		% 	\subfloat[图片14]{\includegraphics[width=0.45\textwidth]{D1-A-3-DataB6.png}\label{fig:image14}}
			
		% 	% 第八行图片
		% 	\subfloat[图片15]{\includegraphics[width=0.45\textwidth]{D1-A-3-DataB7.png}\label{fig:image15}}
		% 	\hfill
		% 	\subfloat[图片16]{\includegraphics[width=0.45\textwidth]{D1-A-3-DataB8.png}\label{fig:image16}}
			
		% 	\caption{16张图片}
		% 	\label{fig:images}
		% \end{figure}

		% \begin{figure}[H]
		% 	\centering
		% 	% 第一行图片
		% 	\subfloat[图片1]{\includegraphics[width=0.45\textwidth]{D1-A-3-Data1.png}\label{fig:image1}}
		% 	\hfill
		% 	\subfloat[图片2]{\includegraphics[width=0.45\textwidth]{D1-A-3-Data2.png}\label{fig:image2}}
			
		% 	% 第二行图片
		% 	\subfloat[图片3]{\includegraphics[width=0.45\textwidth]{D1-A-3-Data3.png}\label{fig:image3}}
		% 	\hfill
		% 	\subfloat[图片4]{\includegraphics[width=0.45\textwidth]{D1-A-3-Data4.png}\label{fig:image4}}
			
		% 	% 第三行图片
		% 	\subfloat[图片5]{\includegraphics[width=0.45\textwidth]{D1-A-3-Data5.png}\label{fig:image5}}
		% 	\hfill
		% 	\subfloat[图片6]{\includegraphics[width=0.45\textwidth]{D1-A-3-Data6.png}\label{fig:image6}}
			
		% 	% 第四行图片
		% 	\subfloat[图片7]{\includegraphics[width=0.45\textwidth]{D1-A-3-Data7.png}\label{fig:image7}}
		% 	\hfill
		% 	\subfloat[图片8]{\includegraphics[width=0.45\textwidth]{D1-A-3-Data8.png}\label{fig:image8}}
			
		% 	\caption{16张图片 - 第1部分}
		% 	\label{fig:images_part1}
		% \end{figure}
		
		% \begin{figure}[H]\ContinuedFloat
		% 	\centering
		% 	% 第五行图片
		% 	\subfloat[图片9]{\includegraphics[width=0.45\textwidth]{D1-A-3-DataB1.png}\label{fig:image9}}
		% 	\hfill
		% 	\subfloat[图片10]{\includegraphics[width=0.45\textwidth]{D1-A-3-DataB2.png}\label{fig:image10}}
			
		% 	% 第六行图片
		% 	\subfloat[图片11]{\includegraphics[width=0.45\textwidth]{D1-A-3-DataB3.png}\label{fig:image11}}
		% 	\hfill
		% 	\subfloat[图片12]{\includegraphics[width=0.45\textwidth]{D1-A-3-DataB4.png}\label{fig:image12}}
			
		% 	% 第七行图片
		% 	\subfloat[图片13]{\includegraphics[width=0.45\textwidth]{D1-A-3-DataB5.png}\label{fig:image13}}
		% 	\hfill
		% 	\subfloat[图片14]{\includegraphics[width=0.45\textwidth]{D1-A-3-DataB6.png}\label{fig:image14}}
			
		% 	% 第八行图片
		% 	\subfloat[图片15]{\includegraphics[width=0.45\textwidth]{D1-A-3-DataB7.png}\label{fig:image15}}
		% 	\hfill
		% 	\subfloat[图片16]{\includegraphics[width=0.45\textwidth]{D1-A-3-DataB8.png}\label{fig:image16}}
			
		% 	\caption{16张图片 - 第2部分}
		% 	\label{fig:images_part2}
		% \end{figure}

		\begin{figure}[H]
			\centering
			\begin{minipage}{0.45\textwidth}
				\subfloat[Data1]{\includegraphics[width=\linewidth]{D1-A-3-Data1.png}\label{fig:image1}}
			\end{minipage}
			\hfill
			\begin{minipage}{0.45\textwidth}
				\subfloat[data2]{\includegraphics[width=\linewidth]{D1-A-3-Data2.png}\label{fig:image2}}
			\end{minipage}
			
			\vspace{1em} % 控制图片上下间距
			
			\begin{minipage}{0.45\textwidth}
				\subfloat[Data3]{\includegraphics[width=\linewidth]{D1-A-3-Data3.png}\label{fig:image3}}
			\end{minipage}
			\hfill
			\begin{minipage}{0.45\textwidth}
				\subfloat[Data4]{\includegraphics[width=\linewidth]{D1-A-3-Data4.png}\label{fig:image4}}
			\end{minipage}
			
			\caption{原始数据质量评定 - 第1部分}
			\label{fig:images_part1}
		\end{figure}
		
		\clearpage % 强制换页
		
		% 第二部分
		\begin{figure}[H]\ContinuedFloat
			\centering
			\begin{minipage}{0.45\textwidth}
				\subfloat[Data5]{\includegraphics[width=\linewidth]{D1-A-3-Data5.png}\label{fig:image5}}
			\end{minipage}
			\hfill
			\begin{minipage}{0.45\textwidth}
				\subfloat[Data6]{\includegraphics[width=\linewidth]{D1-A-3-Data6.png}\label{fig:image6}}
			\end{minipage}
			
			\vspace{1em}
			
			\begin{minipage}{0.45\textwidth}
				\subfloat[Data7]{\includegraphics[width=\linewidth]{D1-A-3-Data7.png}\label{fig:image7}}
			\end{minipage}
			\hfill
			\begin{minipage}{0.45\textwidth}
				\subfloat[Data8]{\includegraphics[width=\linewidth]{D1-A-3-Data8.png}\label{fig:image8}}
			\end{minipage}
			
			\caption{原始数据质量评定 - 第2部分}
			\label{fig:images_part2}
		\end{figure}
		
		\clearpage
		
		% 第三部分
		\begin{figure}[H]\ContinuedFloat
			\centering
			\begin{minipage}{0.45\textwidth}
				\subfloat[DataB1]{\includegraphics[width=\linewidth]{D1-A-3-DataB1.png}\label{fig:image9}}
			\end{minipage}
			\hfill
			\begin{minipage}{0.45\textwidth}
				\subfloat[DataB2]{\includegraphics[width=\linewidth]{D1-A-3-DataB2.png}\label{fig:image10}}
			\end{minipage}
			
			\vspace{1em}
			
			\begin{minipage}{0.45\textwidth}
				\subfloat[DataB3]{\includegraphics[width=\linewidth]{D1-A-3-DataB3.png}\label{fig:image11}}
			\end{minipage}
			\hfill
			\begin{minipage}{0.45\textwidth}
				\subfloat[DataB4]{\includegraphics[width=\linewidth]{D1-A-3-DataB4.png}\label{fig:image12}}
			\end{minipage}
			
			\caption{原始数据质量评定 - 第3部分}
			\label{fig:images_part3}
		\end{figure}
		
		\clearpage
		
		% 第四部分
		\begin{figure}[H]\ContinuedFloat
			\centering
			\begin{minipage}{0.45\textwidth}
				\subfloat[DataB5]{\includegraphics[width=\linewidth]{D1-A-3-DataB5.png}\label{fig:image13}}
			\end{minipage}
			\hfill
			\begin{minipage}{0.45\textwidth}
				\subfloat[DataB6]{\includegraphics[width=\linewidth]{D1-A-3-DataB6.png}\label{fig:image14}}
			\end{minipage}
			
			\vspace{1em}
			
			\begin{minipage}{0.45\textwidth}
				\subfloat[DataB7]{\includegraphics[width=\linewidth]{D1-A-3-DataB7.png}\label{fig:image15}}
			\end{minipage}
			\hfill
			\begin{minipage}{0.45\textwidth}
				\subfloat[DataB8]{\includegraphics[width=\linewidth]{D1-A-3-DataB8.png}\label{fig:image16}}
			\end{minipage}
			
			\caption{原始数据质量评定 - 第4部分}
			\label{fig:images_part4}
		\end{figure}

		通过对每一组数据进行正态分布拟合、X-Y相关性检验，自相关性检验，以及简化卡方值（相关数据均在图上），我们可以确定每一组数据的独立性和随机性。

		可以看到，大部分的数据都具有很好的随机性和独立性，是随机噪声。
		
		但是也有数据并不是保持正态分布，如“DataB1”“DataB5”这两组本底数据的测量。所以，这两组本底以及对应的热噪声测量数据应当删去；Data3 的X 数据列也呈现非正态的分布，但是Y数据列是保持正态分布的，我们可以使用Y数据列进行相关计算。

		这两组数据分别对应了第一组常温的实验数据和第一组低温的实验数据。可能是因为作为常温组和低温组的第一组，未稳定，导致数据分布未保持正态分布。





	\subsection{玻尔兹曼常数的计算}

		在上一步的数据质量评定中，已经计算得到对应的X、Y数据列的标准差，如\cref{tbl:D1-A-2}。

		\begin{table}[htbp]
			\centering
			\begin{tblr}{
			  cells = {c},
			  vline{1-2,4-5} = {-}{},
			  hline{1-2,14} = {-}{},
			}
			File       & X\_std   & Y\_std   & f/Hz     \\
			Data2.csv  & 1.54E-07 & 1.51E-07 & 1145  \\
			Data3.csv  & 1.22E-07 & 1.23E-07 & 11451 \\
			Data4.csv  & 9.38E-08 & 9.57E-08 & 20486 \\
			Data6.csv  & 7.34E-08 & 7.14E-08 & 666   \\
			Data7.csv  & 2.43E-08 & 2.47E-08 & 6666  \\
			Data8.csv  & 3.38E-08 & 3.43E-08 & 36666 \\
			DataB2.csv & 2.30E-08 & 2.35E-08 & 1145  \\
			DataB3.csv & 1.68E-08 & 1.69E-08 & 11451 \\
			DataB4.csv & 1.63E-08 & 1.58E-08 & 20486 \\
			DataB6.csv & 1.56E-08 & 1.57E-08 & 666   \\
			DataB7.csv & 1.57E-08 & 1.55E-08 & 6666  \\
			DataB8.csv & 1.60E-08 & 1.59E-08 & 36666 
			\end{tblr}
			\caption{数据列的标准差}
			\label{tbl:D1-A-2}
		\end{table}

		则根据玻尔兹曼常数的计算公式：
		\[
			k_B = \frac{v_{Nx}^2}{4 \times \text{ENBW} \times T \times R}
		\]

		\begin{itemize}
			\item $v_{Nx}^2$ 的计算：
			
				\[
					v_{Nx}^2 = v_{Rx}^2 [1 + (\omega R_{in} C_{in})^2] - v_{Bx}^2
				\]

				其中$v_{Rx}^2$ 和 $v_{Bx}^2$ 即为我们测量的热噪声的方差（X数据列或者Y数据列均可）；$(\omega R_{in} C_{in})^2$ 项即为输入带宽修正因子，$R_{in} \approx \frac{2}{3} R \ , R = 144 \mathrm{k}\Omega \ , C_{in} = 68 \mathrm{pF}$


			\item ENBW 的计算：
			
				单分量等效噪声带宽（ENBW）与陡降和时间常数 $\tau$ 存在对应关系，实验中我们使用$\tau = 10 \mathrm{ms} \ , Slope = 18 \mathrm{dB/oct}$
				\[
					\Rightarrow \mathrm{ENBW} = \frac{3}{32} / (10 \times 10^{-3}) = 9.375
				\]



		\end{itemize}


		则可以得到玻尔兹曼常数的计算结果，如\cref{tbl:D1-A-3}，$\eta$ 为与标准值的相对误差：

		由于频率为 36666Hz，66666Hz 的数据明显偏大，可能是引入了其他的误差，将其删去，并绘制出玻尔兹曼常数测量值的散点图，如\cref{fig:D1-A-1}

		\begin{table}[htbp]
			\centering
			\begin{tblr}{
			  cells = {c},
			  vline{1-2,8} = {-}{},
			  hline{1,4} = {-}{},
			}
			f/Hz   & 1145     & 11451    & 20486    & 6666     & 66666    & 36666    \\
			$k_B / 10^{23}$  & 1.42 & 1.38 & 1.43 & 1.46 & 2.61 & 1.70 \\
			$\eta$ & 2.551\%  & -0.262\% & 3.665\%  & 5.892\%  & 89.334\% & 23.234\% 
			\end{tblr}
			\caption{玻尔兹曼常数计算结果}
			\label{tbl:D1-A-3}
		\end{table}

		\begin{figure}[htbp]
			\centering
			\includegraphics[width=0.7\textwidth]{D1-A-3-kB-plotting.png}
			\caption{$k_B$ 测量值的散点图}
			\label{fig:D1-A-1}
		\end{figure}

		进行频率与玻尔兹曼常数测量值的相关性分析，有：
		\begin{itemize}
			\item 皮尔逊相关系数：用来衡量 频率 和 $k_B$ 之间的线性关系。
			
				计算得到皮尔逊相关系数为 $-5.4600 \times 10^{-2}$，
			\item 斯皮尔曼相关系数：适用于非线性或不满足正态分布的数据。
				计算得到斯皮尔曼相关系数为 0。
		\end{itemize}

		即，经输入带宽修正后的噪声谱密度符合热噪声的频谱特征，与频率无关。

		若噪声谱密度符合热噪声的频谱特征，与频率无关，则计算玻尔兹曼常数时，可以选择取不同频率下的测量值的平均值。




	\subsection{误差评定}


		\begin{enumerate}
			\item 由误差传递公式：
				\[
					\delta k_B = \sqrt{(\frac{\partial k_B}{\partial T} \delta T)^2 + (\frac{\partial k_B}{\partial R} \delta R)^2}
				\]
				可计算有：
				\begin{align*}
					\frac{\partial k_B}{\partial T} &= -\frac{v_{Nx}^2}{4 \times \text{ENBW} \times T^2 \times R}  	\\
					\frac{\partial k_B}{\partial R} &= -\frac{v_{Nx}^2}{4 \times \text{ENBW} \times T \times R^2}	\\
					% \frac{\partial k_B}{\partial v_{Nx}} &= \frac{2 v_{Nx}}{4 \times \text{ENBW} \times T \times R}
					% k_B = \frac{v_{Nx}^2}{4 \times \text{ENBW} \times T \times R}
				\end{align*}

			\item 考虑测量温度时的B类不确定度：
				\[
					\delta T = \frac{0.1 K}{\sqrt{3}}
				\]
			

			\item 考虑测量电阻时的B类不确定度：
				\[
					\delta R = \frac{0.1 \Omega}{\sqrt{3}}
				\]



			\item 将计算得到的 $\delta k_B$ 绘制为y轴的误差棒添加到图中，如\cref{fig:D1-A-2}
			

				\begin{figure}[htbp]
					\centering
					\includegraphics[width=0.7\textwidth]{D1-A-3-kB-plotting2.png}
					\caption{$k_B$ 测量值的散点图+误差棒}
					\label{fig:D1-A-2}
				\end{figure}
		\end{enumerate}


		在进行数据计算时，我们选择将 X数据列的方差与 Y数据列的方差做平均，作为热噪声的测量值，由于在第一步我们已经进行了数据的独立性检验，以确保所有数据都是独立的，所以这样的操作不会明显降低随机误差。

		





		% \begin{table}[htbp]
		% 	\centering
		% 	\begin{tabular}{|c|cc|} 
		% 	\hline
		% 	File       & X\_std   & Y\_std    \\ 
		% 	\hline
		% 	Data2  & 1.54E-07 & 1.51E-07  \\
		% 	Data3  & 1.22E-07 & 1.23E-07  \\
		% 	Data4  & 9.38E-08 & 9.57E-08  \\
		% 	Data6  & 7.34E-08 & 7.14E-08  \\
		% 	Data7  & 2.43E-08 & 2.47E-08  \\
		% 	Data8  & 3.38E-08 & 3.43E-08  \\
		% 	DataB2 & 2.30E-08 & 2.35E-08  \\
		% 	DataB3 & 1.68E-08 & 1.69E-08  \\
		% 	DataB4 & 1.63E-08 & 1.58E-08  \\
		% 	DataB6 & 1.56E-08 & 1.57E-08  \\
		% 	DataB7 & 1.57E-08 & 1.55E-08  \\
		% 	DataB8 & 1.60E-08 & 1.59E-08  \\
		% 	\hline
		% 	\end{tabular}
		% 	\caption{数据列的标准差}
		% 	\label{tbl:D1-A-2}
		% \end{table}












\clearpage
\subsection{实验报告思考题}


% 思考题1
\begin{question}
	样品盒及传输电缆是否会对本底有贡献？ 用实验现象回答; 处于不同环境（如人的触碰、或插入液氮瓶后） 是否对测量产生影响？
\end{question}

	实验中的热噪声不仅来自电阻本身，还包括系统中的样品盒和传输电缆等部分。这些导电材料也会产生热噪声，叠加在测量信号上。此外，接触电阻可能引入额外噪声，连接不稳或接触不良时影响更大。传输电缆还可能拾取环境中的电磁干扰，如50Hz市电噪声，进一步增加本底噪声。

	本底噪声无需输入带宽修正，因为它主要是由锁相放大器内部产生的噪声，且不受输入电路阻容滤波的影响。而外部信号（如电阻的热噪声）会受到输入阻抗和电容的影响，因此需带宽修正来还原信号强度。

	实验结果也显示，在样品插入液氮瓶后，测量的热噪声显著减小，说明不同环境对测量结果有影响，因为热噪声与温度成正比。






% 思考题2
\begin{question}
	采样时间间隔 $\Delta t$ 约为多少个时间常数时， X意相邻采样读数$X(t)$与 $X(t + \Delta t)$可视为独立无关的？ 如何验证？ 如想缩短测量时间并要求采样读数保持随机噪声特性， 选择时间常数为多长合适？
\end{question}

	$\Delta t$应至少为5-10倍的时间常数 $\tau$ 才能保证相邻采样读数的独立性。

	若要验证采样点的独立性，可以通过计算采样数据的自相关函数来验证。采集一段时间的连续数据，计算数据序列的自相关函数 $R(t)$ 。检查当时间间隔为 $\Delta t$ 时的自相关值是否接近于 0。如果 $R(\Delta t) \approx 0$，则可以认为间隔为 $\Delta t$ 的采样点之间基本独立无关。

	如果希望缩短测量时间，可以适当减少时间常数。选择较小的时间常数可以加快信号的响应速度，使系统在较短时间内达到稳态。

	但需注意，时间常数过小可能导致信号响应变得不稳定或噪声水平增加，因此在实际操作中需要权衡测量精度和时间的要求。若是选择遥测，测量时间则宽松许多，时间常数可以选择放宽至 100 ms，要达到 $10^3$ 以上的数据点，也只需要 20min 左右。



\clearpage
\section{D1 \quad 电阻热噪声及玻尔兹曼常量测量 \quad\heiti 实验总结}

\subsection{实验分工}
	\begin{enumerate}
		\item 戴鹏辉：合作完成实验操作部分。
		\item 杨舒云：合作完成实验操作部分，完成实验方案设计。
	\end{enumerate}


\subsection{实验心得}
	\begin{enumerate}
		\item 这次的实验十分锻炼我们的数据处理以及评定能力。
		
		\item 感谢何老师和柳老师在实验中耐心地解答我们的问题。

	\end{enumerate}

\end{document}
